\begin{codeblock}
struct @\libglobal{tag}@ {};
\end{codeblock}

\begin{itemdescr}
\pnum
\remarks
A defined record's own description.
\end{itemdescr}

\begin{codeblock}
template<class T>
struct @\libglobal{checked}@ {};
\end{codeblock}

\begin{itemdescr}
\pnum
\mandates
\tcode{T} is a type this facility accepts.

\pnum
\remarks
Instantiating \tcode{checked} registers nothing; it only constrains.
\end{itemdescr}

\begin{codeblock}
template<class T>
struct @\libglobal{probe}@ {};
\end{codeblock}

\pnum
A program that instantiates \tcode{probe<T>} is ill-formed unless \tcode{acceptable_v<T>} is \tcode{true}.

\begin{itemdescr}
\pnum
\remarks
The derived paragraph above states the requirement; this says what the type is for.
\end{itemdescr}

\begin{codeblock}
template<class T>
struct @\libglobal{box}@ {
  // \ref{box.mem}, member types
  using type = T;
};
\end{codeblock}

\begin{itemdescr}
\pnum
\remarks
A \tcode{box} is a distinct type for each \tcode{T}; two \tcode{box} types with different \tcode{T} never compare equal.
\end{itemdescr}

\rSec3[box.mem]{Member types}

\indexlibrarymember{type}{box}%
\begin{itemdecl}
using type = T;
\end{itemdecl}

\begin{itemdescr}
\pnum
\remarks
The contained type.
\end{itemdescr}
